

\section{فاصله‌ی واسراشتاین}
متغیر‌های تصادفی 
$X_1, \dots, X_n$
به صورت iid که با توزیع $P$ از مربع
$[0,1]^2$
انتخاب می‌شوند. توزیع تجربی متناظر نقاط را با 
$P_n$
نشان می‌دهیم. فاصله‌ی واسراشتاین این دو توزیع به این صورت تعریف می‌شود:
$$W_1(P_n,P) = \sup_{f: 1-Lipschitz} \mathbb{E}_{P_n}[f] - \mathbb{E}[f]$$
برای سادگی، فرض کنید لیپشیتز بودن نسبت به نرم 
$\infty$
سنجیده می‌شود، یعنی $f$ تابعی ۱-لیپشیتز است اگر و تنها اگر به ازای هر $x$ و $y$ داشته باشیم
$$|f(x) - f(y)| \leq ||x-y||_\infty.$$

قرار دهید
$\mathcal{F}_0 = \{f: 1-Lip, f(0) = 0\}$.

\begin{itemize}
	\item 
	نشان دهید
	$$\mathcal{N}(\mathcal{F}_0, ||.||_\infty, \epsilon) \leq \exp\big(\frac{c}{\epsilon^2}\big).$$
	
	\item 
	نشان دهید که انتگرال دودلی در این مسئله‌ واگراست و کرانی روی فاصله‌ی واسراشتاین نمی‌دهد.
	\item
	با استفاده از تمرین قبل، نشان دهید
	$$\mathbb{E}[W_1(P_n, P)] \leq n^{-1/2} \log(n). $$
\end{itemize}
